\section{Manual de instalación y despliegue} 

    \subsection{Dependencias del sistema}
        El sistema desarrollado depende de varios servicios externos que permiten el funcionamiento del agente conversacional y la comunicación entre los distintos componentes del sistema. 
        Concretamente, el sistema requiere el uso de los siguientes servicios: 
        \begin{itemize}
            \item Discord, utilizado para realizar las llamadas VoIP, capturando y reproduciendo audio de un canal de voz.
            \item Dialogflow, encargado de interpretar la intención del usuario y recopilar la información necesaria para realizar el trámite. 
            \item Google Cloud, necesario para la comunicación entre el bot de Discord y Dialogflow y el uso del servicio de reconocimiento de voz.  
            \item Ngrok, utilizado para exponer el servidor local de desarrollo a Internet, permitiendo la comunicación entre Dialogflow y la plataforma web.
        \end{itemize}
        
        Además, para ejecutar el sistema es necesario disponer del siguiente software: 
        \begin{itemize}
            \item Python 3.13 o superior.
            \item Gestor de paquetes pip.
            \item Navegador web moderno.
            \item Cuenta en Ngrok.
            \item Cuenta en Google Cloud.
            \item Cuenta en Discord.
        \end{itemize}

        Todas las dependencias necesarias para ejecutar el sistema se encuentran definidas en el archivo \textbf{requirements.txt}.

    \subsection{Instalación del proyecto}
        Para instalar el proyecto es necesario clonar el repositorio y preparar el entorno de ejecución de Python. 

        En primer lugar, se debe clonar el repositorio mediante el comando \textbf{git clone \url{https://github.com/humaral/tfg.git}}.
        
        Una vez descargado, se accede al directorio del repositorio con \textbf{cd tfg}. 
        
        Posteriormente, se crea un entorno virtual de Python para aislar las dependencias del proyecto con \textbf{python -m venv venv}.

        Una vez creado el entorno virtual, se activa con \textbf{venv\textbackslash{}Scripts\textbackslash{}activate} en sistemas Windows, y con \textbf{source venv/bin/activate} en Linux o MacOS. 

        Finalmente, se intalan las dependencias necesarias para ejecutar el proyecto con \textbf{pip install -r requirements.txt}.

        Además, es necesario crear un archivo \textbf{.env} que contenga las variables de entorno utilizadas por la aplicación. 
        Debe incluir al menos las siguientes variables:
        \begin{itemize}
            \item \textbf{DIALOGFLOW\_CREDENTIALS}
            \item \textbf{DISCORD\_TOKEN}
            \item \textit{VOSK\_MODEL\_PATH} (\textit{opcional}).
        \end{itemize}

    \subsection{Configuración del agente de Dialogflow}

        El repositorio incluye una copia exportada del agente conversacional utilizado por el sistema, ubicada en \textbf{/scripts\_aux/services/models/Operador.zip}.
        Pasos para importar este agente en Dialogflow se deben seguir los siguientes pasos:
        \begin{enumerate}
            \item Acceder a la \href{https://dialogflow.cloud.google.com}{consola de Dialogflow} - \url{https://dialogflow.cloud.google.com}.
            \item Crear un nuevo agente.
            \item Seleccionar la opción \textbf{Import from ZIP} en la configuración.
            \item Subir el archivo \textbf{Operador.zip}.
        \end{enumerate}

    \subsection{Configuración de Google Cloud}

        Para que la aplicación pueda comunicarse con Dialogflow y para poder usar los servicios de STT es necesario disponer de credenciales de Google Cloud. 
        La configuración de este servicio se consigue siguiendo los siguientes pasos:
        \begin{enumerate}
            \item Acceder a la \href{https://console.cloud.google.com/}{consola de Google Cloud} - \url{https://console.cloud.google.com/}.
            \item Crear un proyecto de Google Cloud.
            \item Activar el servicio de facturación.
            \item Vincular el agente de Dialogflow al proyecto.
            \item Activar el servicio de Speech-to-Text en el proyecto. (\textit{opcional})
            \item Crear una \textbf{Cloud Resource Manager API} asociada a los servicios de Dialogflow y STT y descargar el archivo JSON con sus credenciales.
            \item Añadir en el archivo \textbf{.env} la variable de entorno \textbf{DIALOGFLOW\_CREDENTIALS} con la ruta donde se almacene el archivo JSON con las credenciales.
        \end{enumerate}

    \subsection{Configuración de Vosk}

        En caso de preferir usar un modelo local para realizar el Speech-to-Text en vez del servicio de Google, se puede conseguir siguiendo estos pasos:
        \begin{enumerate}
            \item Descargar el \href{https://alphacephei.com/vosk/models}{modelo Vosk} en \url{https://alphacephei.com/vosk/models}.
            \item Descomprimir la carpeta y añadirla al proyecto.
            \item Añadir en el archivo \textbf{.env} la variable de entorno \textbf{VOSK\_MODEL\_PATH} con la ruta donde se almacene la carpeta con el modelo.
            \item Acceder al archivo \textbf{scripts\_aux\textbackslash{}services\textbackslash{}discord\_bot.py} y cambiar la variable global \textbf{USE\_LOCAL\_STT=True}.
        \end{enumerate}

    \subsection{Configuración del bot de Discord}

        El código del bot se encuentra en \textbf{scripts\_aux\textbackslash{}services\textbackslash{}discord\_bot.py} y para utilizarlo hay que seguir los siguientes pasos:
        \begin{enumerate}
            \item Acceder a \href{https://discord.com/developers}{Discord Developer Portal} - \url{https://discord.com/developers}.
            \item Crear una nueva aplicación.
            \item Añadir un bot a la aplicación.
            \item Copiar el token de autenticación del bot y añadirlo en \textbf{.env} en la variable \textbf{DISCORD\_TOKEN}.
            \item Crear un servidor de Discord.
            \item Generar un enlace de invitación desde el portal de desarrolladores e invitar al bot al servidor de Discord.
            \item Crear un canal de voz en el servidor de Discord, limitado a 2 usuarios.
        \end{enumerate}
        Para que el bot pueda funcionar correctamente, deberá tener permisos en el servidor de Discord para conectarse a canales de voz, hablar y mover miembros.

    \subsection{Configuración de Ngrok}

        Antes de iniciar Ngrok por primera vez, hay que realizar los siguientes pasos: 
        \begin{enumerate}
            \item Acceder al \href{https://dashboard.ngrok.com/}{dashboard de Ngrok} - \url{https://dashboard.ngrok.com/}.
            \item Copiar el Authtoken y añadirlo al entorno con el comando \textbf{ngrok config add-authtoken TU\_TOKEN\_NGROK}.
            \item Ir a la pestaña Domains en el dashboard de Ngrok y copiar la URL pública donde se levantará la aplicación web.
            \item En la consola de Dialogflow, acceder a la pestaña de fulfillment y añadir la URL del webhook de la aplicación web - \textbf{https://TU\_URL\_NGROK/webhook}.
        \end{enumerate}

    \subsection{Despliegue del sistema}
        Para ejecutar el sistema es necesario iniciar los distintos componentes del sistema en orden. 
        
        En primer lugar, hay que iniciar el servidor de correo SMTPD utilizado para simular el envío de correos electrónicos con \textbf{py -m aiosmtpd -n -l localhost:1025}, recibiéndolos en texto plano en el terminal donde se ejecuta:

        Posteriormente, se inicia el servidor Flask con \textbf{py app.py}, quedando la aplicación web disponible para acceder desde un navegador en la dirección \textbf{http://localhost:5000}.

        Una vez iniciada la aplicación web, se debe exponer el servidor local a Internet con el servicio de Ngrok, ejecutando \textbf{ngrok http 5000}.

        Por último, se inicializa el bot de Discord ejecutando \textbf{py .\textbackslash{}scripts\_aux\textbackslash{}services\textbackslash{}discord\_bot.py}.
        El bot se conectará al servidor de Discord y cuando el usuario entre al canal de voz, dará comienzo la llamada.

\section{Manual de usuario}
